%%% BEGIN fig:intro-dependency
\begin{figure}[htbp]
\centering
% Layout generated by agent figA from FINAL_OUTLINE_v2: node fill = rigour tag; grey arrows = transitive
% reduction of the \uses relation among the displayed results; circled digits = tools of column 1 used;
% orange dashed = numbered assumptions (premises of R_cond theorems); thick red = bottleneck assumptions.
\begin{tikzpicture}[
  fa res/.style={draw=#1!75!black,fill=#1!14,line width=0.6pt,rounded corners=2.5pt,
    text width=1.74cm,align=flush center,inner xsep=1pt,inner ysep=2pt,font=\fontsize{7.5}{8.8}\selectfont},
  fa cond/.style={dash pattern=on 2.2pt off 1.4pt},
  fa asm/.style={draw=orange!75!black,fill=orange!7,line width=0.6pt,rounded corners=2pt,
    dash pattern=on 2.2pt off 1.4pt,align=flush center,inner xsep=2.5pt,inner ysep=2pt,font=\scriptsize},
  fa dep/.style={draw=gray!55!black,line width=0.5pt,rounded corners=2.5pt,-{Stealth[length=3.4pt,width=3pt]}},
  fa bot/.style={draw=red!70!black,line width=1.2pt,rounded corners=3pt,-{Stealth[length=4.6pt,width=4.2pt]}},
  fa use/.style={draw=orange!75!black,line width=0.8pt,dash pattern=on 2.5pt off 1.5pt,rounded corners=3pt,
    -{Stealth[length=4pt,width=3.6pt]}},
  fa tool/.style={inner sep=0pt,font=\small},
  fa head/.style={align=center,font=\scriptsize,inner sep=1pt}]
  \node[fa head] at (0.92,1.42) {\textbf{工具}\\[1pt]{\tiny 第 2--3 章}};
  \node[fa head] at (3.08,1.42) {\textbf{红外结构}\\[1pt]{\tiny 第 4--5 章}};
  \node[fa head] at (5.47,1.42) {\textbf{严格核心}\\[1pt]{\tiny 第 6--8 章}};
  \node[fa head] at (7.86,1.42) {\textbf{因子化机器}\\[1pt]{\tiny 第 9 章与附录 N、O}};
  \node[fa head] at (10.32,1.42) {\textbf{过程}\\[1pt]{\tiny 第 10--15 章}};
  \node[fa head] at (12.71,1.42) {\textbf{边界}\\[1pt]{\tiny 第 18 章}};
  \node[fa head] at (15.05,1.42) {\textbf{前沿}\\[1pt]{\tiny 第 22--24 章}};
  \node[fa res=OliveGreen] (ST) at (0.92,0.00) {ST 恒等式\\[-0.4pt]\mbox{{\scriptsize{\fontspec{FandolSong-Regular.otf}①}}\,\tiny\color{gray!25!black}\cref{thm:prelim-slavnov-taylor}}};
  \node[fa res=OliveGreen] (BETA) at (0.92,-1.15) {$\beta$ 函数\\[-0.4pt]\mbox{{\scriptsize{\fontspec{FandolSong-Regular.otf}②}}\,\tiny\color{gray!25!black}\cref{prop:prelim-beta}}};
  \node[fa res=OliveGreen] (WIL) at (0.92,-2.42) {Wilson 线的\\动量空间规则\\[-0.4pt]\mbox{{\scriptsize{\fontspec{FandolSong-Regular.otf}③}}\,\tiny\color{gray!25!black}\cref{prop:prelim-wilson-rules}}};
  \node[fa res=OliveGreen] (CUT) at (0.92,-3.70) {切割规则\\[-0.4pt]\mbox{{\scriptsize{\fontspec{FandolSong-Regular.otf}④}}\,\tiny\color{gray!25!black}\cref{thm:prelim-cutting-rules}}};
  \node[fa res=OliveGreen] (COMP) at (0.92,-4.97) {复合算符\\重整化\\[-0.4pt]\mbox{{\scriptsize{\fontspec{FandolSong-Regular.otf}⑤}}\,\tiny\color{gray!25!black}\cref{thm:prelim-composite-renorm}}};
  \node[fa res=OliveGreen] (PC) at (3.08,0.00) {红外幂次计数\\[-0.4pt]\mbox{\tiny\color{gray!25!black}\cref{prop:ir-power-counting}}};
  \node[fa res=OliveGreen] (LAN) at (3.08,-1.15) {Landau 方程\\[-0.4pt]\mbox{\tiny\color{gray!25!black}\cref{thm:ir-landau-momentum}}};
  \node[fa res=RoyalBlue] (LR) at (3.08,-2.30) {领头区域\\[-0.4pt]\mbox{\tiny\color{gray!25!black}\cref{thm:ir-leading-regions}\,{\fontsize{5.5}{6}\selectfont{\fontspec{FandolSong-Regular.otf}①}}}};
  \node[fa res=RoyalBlue] (SUD) at (3.08,-3.56) {Sudakov\\形状因子\\[-0.4pt]\mbox{\tiny\color{gray!25!black}\cref{thm:sudakov-factorization}\,{\fontsize{5.5}{6}\selectfont{\fontspec{FandolSong-Regular.otf}③}}}};
  \node[fa res=OliveGreen] (PML) at (5.47,0.00) {PDF 的矩\\[-0.4pt]\mbox{\tiny\color{gray!25!black}\cref{thm:pdf-moments-local}\,{\fontsize{5.5}{6}\selectfont{\fontspec{FandolSong-Regular.otf}③⑤}}}};
  \node[fa res=OliveGreen] (DGL) at (5.47,-1.20) {DGLAP 方程\\[-0.4pt]\mbox{\tiny\color{gray!25!black}\cref{thm:pdf-dglap}\,{\fontsize{5.5}{6}\selectfont{\fontspec{FandolSong-Regular.otf}①②}}}};
  \node[fa res=OliveGreen] (EE) at (5.47,-2.30) {$R(s)$ 有限性\\[-0.4pt]\mbox{\tiny\color{gray!25!black}\cref{thm:ee-R-finite}\,{\fontsize{5.5}{6}\selectfont{\fontspec{FandolSong-Regular.otf}①②③}}}};
  \node[fa res=RoyalBlue] (OPE) at (5.47,-3.45) {欧氏 OPE\\[-0.4pt]\mbox{\tiny\color{gray!25!black}\cref{thm:dis-ope}\,{\fontsize{5.5}{6}\selectfont{\fontspec{FandolSong-Regular.otf}①⑤}}}};
  \node[fa res=OliveGreen] (MOM) at (5.47,-4.60) {DIS 矩因子化\\[-0.4pt]\mbox{\tiny\color{gray!25!black}\cref{thm:dis-moments}\,{\fontsize{5.5}{6}\selectfont{\fontspec{FandolSong-Regular.otf}③}}}};
  \node[fa res=OliveGreen,fa cond] (MNP) at (5.47,-5.86) {DIS 矩的\\非微扰形式\\[-0.4pt]\mbox{\tiny\color{gray!25!black}\cref{thm:dis-moments-np}\,{\fontsize{5.5}{6}\selectfont{\fontspec{FandolSong-Regular.otf}①②⑤}}}};
  \node[fa res=OliveGreen,fa cond] (SUBN) at (7.86,0.00) {条件减除定理\\[-0.4pt]\mbox{\tiny\color{gray!25!black}\cref{thm:appN-subtraction}}};
  \node[fa res=RoyalBlue] (SUB) at (7.86,-1.20) {减除定理\\[-0.4pt]\mbox{\tiny\color{gray!25!black}\cref{thm:method-subtraction}\,{\fontsize{5.5}{6}\selectfont{\fontspec{FandolSong-Regular.otf}⑤}}}};
  \node[fa res=RoyalBlue] (TPL) at (7.86,-2.30) {因子化模板\\[-0.4pt]\mbox{\tiny\color{gray!25!black}\cref{thm:method-template}}};
  \node[fa res=RoyalBlue] (GDEF) at (7.86,-3.45) {围道变形判据\\[-0.4pt]\mbox{\tiny\color{gray!25!black}\cref{prop:method-glauber-deformation}\,{\fontsize{5.5}{6}\selectfont{\fontspec{FandolSong-Regular.otf}④}}}};
  \node[fa res=OliveGreen,fa cond] (GLO) at (7.86,-4.60) {Glauber 相消\\[-0.4pt]\mbox{\tiny\color{gray!25!black}\cref{thm:appO-main}}};
  \node[fa res=RoyalBlue] (WARD) at (7.86,-5.75) {Ward 恒等式\\[-0.4pt]\mbox{\tiny\color{gray!25!black}\cref{lem:method-ward-multi}\,{\fontsize{5.5}{6}\selectfont{\fontspec{FandolSong-Regular.otf}①③}}}};
  \node[fa res=RoyalBlue] (DIS) at (10.32,-1.20) {DIS 因子化\\[-0.4pt]\mbox{\tiny\color{gray!25!black}\cref{thm:disall-factorization}}};
  \node[fa res=RoyalBlue] (SIA) at (10.32,-2.30) {$e^+e^-\to hX$\\[-0.4pt]\mbox{\tiny\color{gray!25!black}\cref{thm:coll-sia}}};
  \node[fa res=RoyalBlue] (DY) at (10.32,-3.45) {单举 DY\\[-0.4pt]\mbox{\tiny\color{gray!25!black}\cref{thm:dy-factorization}}};
  \node[fa res=RoyalBlue] (TMD) at (10.32,-4.60) {TMD 因子化\\[-0.4pt]\mbox{\tiny\color{gray!25!black}\cref{thm:tmd-factorization}}};
  \node[fa res=RoyalBlue] (THR) at (10.32,-5.75) {推力分布\\[-0.4pt]\mbox{\tiny\color{gray!25!black}\cref{thm:scet-thrust}}};
  \node[fa res=BrickRed] (BTMD) at (12.71,-2.58) {$hh\to HHX$\\TMD 因子化\\[-0.4pt]\mbox{\tiny\color{gray!25!black}\cref{thm:break-tmd}\,{\fontsize{5.5}{6}\selectfont{\fontspec{FandolSong-Regular.otf}③}}}};
  \node[fa res=Plum] (COL) at (12.71,-3.83) {有色末态\\共线因子化\\[-0.4pt]\mbox{\tiny\color{gray!25!black}\cref{conj:dy-coloured-final}}};
  \node[fa res=BrickRed] (BSP) at (12.71,-5.12) {类空严格\\共线因子化\\[-0.4pt]\mbox{\tiny\color{gray!25!black}\cref{prop:break-spacelike}}};
  \node[fa res=RoyalBlue] (PIN) at (15.05,0.00) {物理区\\钳夹判据\\[-0.4pt]\mbox{\tiny\color{gray!25!black}\cref{thm:front-collins-pinch}}};
  \node[fa res=RoyalBlue] (QTMD) at (15.05,-1.36) {准 TMD\\因子化\\[-0.4pt]\mbox{\tiny\color{gray!25!black}\cref{thm:lamet-quasi-tmd}\,{\fontsize{5.5}{6}\selectfont{\fontspec{FandolSong-Regular.otf}③⑤}}}};
  \node[fa res=RoyalBlue] (LAM) at (15.05,-2.58) {LaMET 匹配\\[-0.4pt]\mbox{\tiny\color{gray!25!black}\cref{thm:lattice-lamet}}};
  \node[fa res=RoyalBlue] (TO) at (15.05,-3.83) {软因子\\时序无关性\\[-0.4pt]\mbox{\tiny\color{gray!25!black}\cref{thm:yzliu-time-ordering}}};
  \begin{scope}[on background layer]
    \draw[fa dep] (WIL.south) -- (CUT.north);
    \draw[fa dep] (LAN.south) -- (LR.north);
    \draw[fa dep] (LR.south) -- (SUD.north);
    \draw[fa dep] (PML.south) -- (DGL.north);
    \draw[fa dep] (OPE.south) -- (MOM.north);
    \draw[fa dep] (SUBN.south) -- (SUB.north);
    \draw[fa dep] (SUB.south) -- (TPL.north);
    \draw[fa dep] (DY.south) -- (TMD.north);
    \draw[fa dep] (MNP.north) -- (MOM.south);
    \draw[fa dep] (GDEF.north) -- (TPL.south);
    \draw[fa dep] (GLO.north) -- (GDEF.south);
    \draw[fa dep] (WARD.north) -- (GLO.south);
    \draw[fa dep] (PC.west) -| (2.000,-1.000) |- (LR.west);
    \draw[fa dep] (TPL.east) -| (8.900,-1.700) |- (DIS.west);
    \draw[fa dep] (TPL.east) -- (SIA.west);
    \draw[fa dep] (TPL.east) -| (8.900,-2.900) |- ($(DY.west)+(0,0.11)$);
    \draw[fa dep] ($(GDEF.east)+(0,-0.11)$) -| (9.060,-4.700) |- (THR.west);
    \draw[fa dep] ($(DY.east)+(0,0.11)$) -| (11.500,-3.600) |- (COL.west);
    \draw[fa dep] (DGL.east) -| (6.500,-1.750) |- ($(TPL.west)+(0,0.11)$);
    \draw[white,line width=2pt] (4.530,-0.530) -- (6.590,-0.530);
    \draw[white,line width=2pt] (6.830,-0.530) -- (13.840,-0.530);
    \draw[white,line width=2pt] (6.620,-0.720) -- (13.650,-0.720);
    \draw[white,line width=2pt] (4.530,-1.760) -- (12.560,-1.760);
    \draw[white,line width=2pt] (4.530,-3.940) -- (6.680,-3.940);
    \draw[white,line width=2pt] (4.260,-4.110) -- (11.520,-4.110);
    \draw[fa dep] (LR.east) -| (4.410,-1.000) |- (6.710,-0.530) |- (SUBN.west);
    \draw[fa dep] (LR.east) -| (4.410,-1.000) |- (13.960,-0.530) |- (QTMD.west);
    \draw[fa dep] (LR.east) -| (4.410,-1.500) |- (12.710,-1.760) -| (BTMD.north);
    \draw[fa dep] (LR.east) -| (4.410,-3.400) |- (6.800,-3.940) |- ($(GLO.west)+(0,0.11)$);
    \draw[fa dep] (DGL.east) -| (6.500,-0.900) |- (13.770,-0.720) |- (LAM.west);
    \draw[fa dep] (SUD.east) -| (4.140,-3.900) |- (11.640,-4.110) |- (BSP.west);
    \draw[fa dep] (LAN.east) -| (4.270,0.820) -| (PIN.north);
  \end{scope}
  \node[fa asm] (A1) at (3.05,-7.42) {非微扰 OPE [C]\\[-0.4pt]{\tiny\color{gray!25!black}\cref{asm:intro-np-ope}}};
  \node[fa asm] (A2) at (6.45,-7.42) {Glauber 区 [A]\\[-0.4pt]{\tiny\color{gray!25!black}\cref{asm:appO-mixing,asm:appO-deformation,asm:appO-rewritten-cs}}};
  \node[fa asm] (A3) at (10.32,-7.42) {桥接假设 [C]\\[-0.4pt]{\tiny\color{gray!25!black}\cref{asm:intro-bridge}}};
  \node[fa asm] (A4) at (12.85,-7.42) {Collins 条件 [A]\\[-0.4pt]{\tiny\color{gray!25!black}\cref{asm:appN-collins}}};
  \begin{scope}[on background layer]
    \draw[fa use] (A1.north) -- (MNP.south);
    \draw[fa use] (A2.north -| 6.61,0) |- ($(GLO.west)+(0,-0.11)$);
    \draw[fa bot] (A2.east) -| (9.230,-5.000) |- ($(DY.west)+(0,-0.11)$);
    \draw[fa bot] (A2.east) -| (9.230,-5.000) |- (TMD.west);
    \draw[fa bot] (A3.north) -- (THR.south);
    \draw[fa bot] (A4.west) -| (11.380,-6.000) |- (DIS.east);
    \draw[fa bot] (A4.west) -| (11.380,-6.000) |- (SIA.east);
    \draw[fa bot] (A4.west) -| (11.380,-6.000) |- ($(DY.east)+(0,-0.11)$);
    \draw[fa bot] (A4.west) -| (11.380,-6.000) |- (TMD.east);
  \end{scope}
  \begin{scope}[shift={(14.15,-4.62)},font=\fontsize{7.5}{9}\selectfont]
    \filldraw[fill=OliveGreen!14,draw=OliveGreen!75!black,line width=0.5pt,rounded corners=1pt] (0,-0.095) rectangle ++(0.36,0.19);
    \node[anchor=west,inner sep=1pt] at (0.42,0.000) {R 严格};
    \filldraw[fill=RoyalBlue!14,draw=RoyalBlue!75!black,line width=0.5pt,rounded corners=1pt] (0,-0.445) rectangle ++(0.36,0.19);
    \node[anchor=west,inner sep=1pt] at (0.42,-0.350) {A 全阶论证};
    \filldraw[fill=BurntOrange!14,draw=BurntOrange!75!black,line width=0.5pt,rounded corners=1pt] (0,-0.795) rectangle ++(0.36,0.19);
    \node[anchor=west,inner sep=1pt] at (0.42,-0.700) {P 部分};
    \filldraw[fill=Plum!14,draw=Plum!75!black,line width=0.5pt,rounded corners=1pt] (0,-1.145) rectangle ++(0.36,0.19);
    \node[anchor=west,inner sep=1pt] at (0.42,-1.050) {C 猜想};
    \filldraw[fill=BrickRed!14,draw=BrickRed!75!black,line width=0.5pt,rounded corners=1pt] (0,-1.495) rectangle ++(0.36,0.19);
    \node[anchor=west,inner sep=1pt] at (0.42,-1.400) {V 破坏};
    \draw[gray!60!black,line width=0.5pt,dash pattern=on 2.2pt off 1.4pt,rounded corners=1pt] (0,-1.845) rectangle ++(0.36,0.19);
    \node[anchor=west,inner sep=1pt] at (0.42,-1.75) {条件定理};
    \node[anchor=west,inner sep=1pt] at (-0.03,-2.1) {{\fontspec{FandolSong-Regular.otf}①⑤}\ 所用工具};
    \draw[fa bot] (0,-2.45) -- (0.36,-2.45);
    \node[anchor=west,inner sep=1pt] at (0.42,-2.45) {瓶颈假设};
    \draw[fa use] (0,-2.8) -- (0.36,-2.8);
    \node[anchor=west,inner sep=1pt] at (0.42,-2.8) {前提};
  \end{scope}
\end{tikzpicture}
\caption{本书主要结论的逻辑依赖。框的颜色表示严格性标签（绿 R、蓝 A、橙 P、紫 C、红 V），虚线框为 $R_{\rm cond}$ 条件定理；灰色箭头是各条目依赖记录在所示结论之间的传递约化，引用编号后的圈码表示用到左列的哪几项工具（ST 为 Slavnov--Taylor 恒等式）。下方橙色虚线框为编号假设，其中“Glauber 区”框代表 Glauber--软混合、围道形变与改写形式的 Collins 条件三条假设；橙色虚线箭头指向以之为前提的条件定理，粗红线标出决定各全阶论证地位的瓶颈假设。完整的依赖图见附录~P。}
\label{fig:intro-dependency}
\end{figure}
%%% END fig:intro-dependency
